\section{Методика экспериментов}
\label{sec:methodology}

\subsection{Объект сравнения и контрольные вычисления}

Цель эксперимента состоит в проверке двух свойств системы: корректности
конфиденциального вывода по заданному дереву и зависимости затрат на вывод
от размера дерева и пакета запросов. Единицей сравнения служит прогноз для
одного объекта. Во всех вариантах используются одна версия дерева, один
порядок признаков и одно правило перехода в левое поддерево при равенстве
признака и порога. Обучение завершается до начала измерений.

Для объекта $x_i$ регистрируются три метки класса. Через $y_i^{(0)}$
обозначается результат открытого обхода исходного дерева, через $y_i^{(q)}$
--- результат открытого обхода после кодирования признаков и порогов
по формуле~\eqref{eq:fixed-point}, через $y_i^{(s)}$~--- восстановленная
метка конфиденциального вычисления. В целочисленном эталоне применяется
обычный последовательный обход по выбранной ветви. Это позволяет проверить
вычисление всех узлов и скрытый выбор листа по независимой реализации
той же функции.

На наборе из $n$ объектов вычисляются две доли расхождений:
\[
  e_{\mathrm{impl}}=\frac{1}{n}\sum_{i=1}^{n}
       [y_i^{(s)}\ne y_i^{(q)}], \qquad
  e_{\mathrm{quant}}=\frac{1}{n}\sum_{i=1}^{n}
       [y_i^{(q)}\ne y_i^{(0)}].
\]
Квадратные скобки обозначают индикатор события. Критерий корректности
реализации на контрольном наборе имеет вид $e_{\mathrm{impl}}=0$.
Величина $e_{\mathrm{quant}}$ характеризует влияние представления чисел
на итоговый прогноз. При наличии истинных меток дополнительно определяется
точность классификации каждого варианта по отношению к разметке. Совпадение
двух реализаций и правильность прогноза относительно разметки отвечают
на разные экспериментальные вопросы.

\subsection{Формирование контрольных наборов}

Для проверки сравнений формируется синтетический набор с заранее
известными путями от корня до листьев. Для каждого проверяемого порога
выбираются значения по обе стороны от него, точное равенство и пара
различных исходных значений с одинаковым кодом. Включаются отрицательные
признаки и пороги. Для проверки внутреннего узла остальные признаки
подбираются так, чтобы открытый обход достигал этого узла. В тестовом
дереве листьям присваиваются различные коды: тогда ошибка выбора пути
проявляется в итоговой метке. Отдельный тест с повторяющимися классами
проверяет, что получателю выдаётся код класса, предусмотренный моделью.

Для исследования на размеченных данных предусматривается табличная
выборка с числовыми признаками. До обучения фиксируются версия набора,
правила обработки пропусков, разбиение на обучающую и контрольную части
и начальное состояние генератора, используемого при разбиении.
Преобразования признаков определяются по обучающей части и затем
применяются к контрольной. Для каждого обученного дерева сохраняются
топология, индексы признаков, пороги и метки листьев. Это обеспечивает
повторное использование одной модели во всех трёх вычислениях.

\subsection{План измерений}

Варьируются глубина дерева $d$, фактическое число внутренних узлов $m$,
размер пакета $B$ и число дробных разрядов $f$. Для оценки влияния
округления изменяют $f$ при фиксированных дереве и входах. Для оценки
производительности фиксируют $f$ и повторяют вывод на одном и том же
пакете с различными конфигурациями дерева. Наряду с глубиной указывается
$m$, поскольку деревья одинаковой глубины могут требовать разного числа
сравнений. До начала серии проверяется представимость значений и разностей,
используемых для знакового сравнения.

Время подготовки и загрузки модели учитывается отдельно. Задержка пакета
измеряется от начала формирования долей подготовленных входов до
восстановления последней метки у получателя. При длительности обработки
$t_B$ пропускная способность равна $B/t_B$. Первый запуск регистрируется
отдельно; после прогрева выполняются повторные измерения при фиксированных
версиях программных компонентов и одинаковом выделении ресурсов.
Для каждой конфигурации приводятся число повторений, медиана задержки
и межквартильный размах.

Для каждого вычислительного процесса регистрируются процессорное время
и пиковое потребление оперативной памяти. При наличии инструментированного
межсерверного транспорта объём обмена определяется суммой отправленных
байтов по трём серверам, без повторного учёта при приёме. Число раундов
определяется по последовательным этапам протокольного взаимодействия;
несколько сообщений одного этапа могут передаваться параллельно.
Если локальный режим исполняет операции в общей среде без наблюдаемого
сетевого обмена, сетевые показатели для такой конфигурации обозначаются
как недоступные. Переход к измерению производительности предусмотрен
после прохождения функциональных проверок раздела~\ref{sec:trial}.


Для первой серии выбираются полные бинарные деревья глубины
$d\in\{1,2,4,6\}$ и пакеты $B\in\{1,16,64\}$. Декартово произведение
этих множеств задаёт 12 конфигураций; значения $m$ равны соответственно
1, 3, 15 и 63. Во всех конфигурациях используется единое $f$,
прошедшее проверку диапазона. Для каждой конфигурации планируется один
прогревочный запуск и 20 измеряемых повторов на одном фиксированном пакете.
В отдельной серии на фиксированных дереве и наборе объектов изменяется
только $f$; выбор допустимых значений определяется интервалами признаков
и порогов.
